\section{\method{}-Guided Trap-Aware Recovery}
\label{sec:trap_intervention}

The \method{} analysis asks where high- and low-outcome traffic concentrates on a shared task map.
We next test whether that map can guide a small recovery pipeline: act only when a live rollout reaches a state resembling historical traps, and use only information already visible in the agent's context.
Because \S\ref{sec:rq3} shows that \swe{} most strongly demands recovery after trap exposure, we instantiate trap-aware recovery on \swe{} and measure official resolved rate after trap-triggered continuations.
The goal is not a universal repair algorithm, but a test of whether graph-derived traps identify states where a lightweight, non-oracular continuation change helps.

\begin{table*}[!htbp]
\centering
\small
\setlength{\tabcolsep}{5pt}
\renewcommand{\arraystretch}{1.15}
\begin{tabular}{lc|c|cc}
\toprule
\textbf{Provider} & $n_{\text{fired}}$ & \textbf{Baseline} & \textbf{Hot} & \textbf{Note} \\
 & & {\scriptsize($T{=}0.6$, no note)} & {\scriptsize($T{=}0.9$, no note)} & {\scriptsize($T{=}0.6$, note)} \\
\midrule
\multicolumn{5}{c}{\emph{Per-provider fired subset (each provider's own complete set)}} \\
\midrule
Qwen3.6-35B-A3B & 214 & 29.4 & 30.8 {\scriptsize\textcolor{gaincolor}{(+1.4)}} & \cellcolor{bestcolor}\textbf{34.1} {\scriptsize\textcolor{gaincolor}{(+4.7)}} \\
GLM-5.1 & 196 & 48.0 & \cellcolor{bestcolor}\textbf{52.6} {\scriptsize\textcolor{gaincolor}{(+4.6)}} & 49.5 {\scriptsize\textcolor{gaincolor}{(+1.5)}} \\
DeepSeek-V4-Pro & 206 & 44.7 & \cellcolor{bestcolor}\textbf{48.1} {\scriptsize\textcolor{gaincolor}{(+3.4)}} & 45.6 {\scriptsize\textcolor{gaincolor}{(+1.0)}} \\
\midrule
\rowcolor{highlight}
Pooled (3-prov, $n{=}616$ cells) & 616 & 40.4 & \textbf{43.5} {\scriptsize\textcolor{gaincolor}{(+3.1)}} & 42.9 {\scriptsize\textcolor{gaincolor}{(+2.4)}} \\
\midrule
\multicolumn{5}{c}{\emph{Common-fired subset (same 96 instances, all three providers)}} \\
\midrule
Qwen3.6-35B-A3B & 96 & 31.2 & 31.2 \phantom{{\scriptsize{(+0.0)}}} & \cellcolor{bestcolor}\textbf{34.4} {\scriptsize\textcolor{gaincolor}{(+3.1)}} \\
GLM-5.1 & 96 & 50.0 & \cellcolor{bestcolor}\textbf{56.2} {\scriptsize\textcolor{gaincolor}{(+6.2)}} & 52.1 {\scriptsize\textcolor{gaincolor}{(+2.1)}} \\
DeepSeek-V4-Pro & 96 & 41.7 & \cellcolor{bestcolor}\textbf{46.9} {\scriptsize\textcolor{gaincolor}{(+5.2)}} & 44.8 {\scriptsize\textcolor{gaincolor}{(+3.1)}} \\
\midrule
\rowcolor{highlight}
Pooled (3-prov, $n{=}288$ cells) & 288 & 41.0 & \textbf{44.8} {\scriptsize\textcolor{gaincolor}{(+3.8)}} & 43.8 {\scriptsize\textcolor{gaincolor}{(+2.8)}} \\
\bottomrule
\end{tabular}
\caption{\method{}-guided trap-aware recovery on the full 500-instance \swe{} Verified pool. Cells report official resolved rate (\%); parentheses show paired gain over Baseline. Gains are computed from unrounded paired outcomes, so they may differ by 0.1 pp from subtracting displayed rounded rates. Best non-Baseline cells are bold-shaded. Appendix~\ref{app:trap_intervention_more} reports uncertainty; Appendix~\ref{app:trap_case_studies} gives cases.}
\label{tab:swe_4arm_intervention}
\end{table*}

\subsection{From Trap Regions to Runtime Triggers}
\label{sec:intervention_design}

\method{} converts the trap overlay into a runtime detector by storing canonicalized key sets from trap-side states, such as \texttt{ACTION:edit}, \texttt{OBS:error}, \texttt{TOOL:<name>}, and local file-path cues.
Generic phase and non-tool keys are stripped so the detector matches recurring failure contexts rather than full traces.
A second library comes from high-committor non-trap states.
At a live step $s$, the detector computes the candidate key set $q_s$ and scores
\begin{align}
\mathrm{sim}_{\mathrm{trap}}(s) &= \max_{u\in\mathcal{T}} J_{\idf}(q_s,u), \nonumber\\
\mathrm{sim}_{\mathrm{core}}(s) &= \max_{u\in\mathcal{N}} J_{\idf}(q_s,u),
\end{align}
where $\mathcal{T}$ and $\mathcal{N}$ are the trap and non-trap libraries.
A trigger fires when trap similarity clears a fixed threshold and exceeds core similarity by a margin (Appendix Table~\ref{tab:detector_hparams}).
For \swe{}, the current step must also carry a local file-path cue, and the candidate action must be an edit or submit action.
These local checks keep the trigger tied to observable graph signatures while reducing fires on generic exploration.

\subsection{Prefix-Fork Recovery Policy}
\label{sec:prefix_fork}

When the detector fires, we snapshot the docker workspace and conversation prefix, then continue from that identical state under three policies.
The no-intervention \emph{Baseline} continues at $T{=}0.6$ with no added note.
\emph{Hot} raises temperature to $T{=}0.9$ with top-$p{=}0.95$.
\emph{Note} keeps $T{=}0.6$ and appends a conservative diagnosis note drawn only from the agent's recent log: the last action, observation signal, implicated file or target keys, detector confidence, and a family-level diagnosis.
The note asks the agent to re-read the failing evidence, localize narrowly, make one minimal change, and revise or discard the edit if the evidence does not support it.
No arm receives a gold patch, verified-test list, or oracle information.

This prefix-fork design compares the same trap state with no pipeline action, extra sampling diversity, or evidence-grounded localization.
We report official \swe{} resolved rate in Table~\ref{tab:swe_4arm_intervention}; Appendix~\ref{app:trap_intervention_more} provides paired uncertainty estimates.

\subsection{Downstream Accuracy Improves at Trap-Triggered States}
\label{sec:swe_intervention}

The experiment uses the full 500-instance \swe{} Verified pool at seed $11$, spanning 12 repositories and three continuation providers not used in the observational graph analysis: Qwen3.6-35B-A3B, GLM-5.1, and DeepSeek-V4-Pro.
Table~\ref{tab:swe_4arm_intervention} reports each provider's complete fired subset and a common-fired subset of 96 instances where all providers fire and complete all three arms.

\method{}-triggered recovery improves downstream resolution on fired states.
In the per-provider fired subset, pooled Hot raises resolved rate from $40.4\%$ to $43.5\%$ ($+3.1$ pp, $p{=}0.016$), and pooled Note reaches $42.9\%$ ($+2.4$ pp, $p{=}0.064$).
The common-fired subset shows the same pattern: Hot increases pooled resolved rate from $41.0\%$ to $44.8\%$ ($+3.8$ pp, $p{=}0.045$).
Thus a trap region discovered from shared offline trajectories can become a runtime trigger that improves a downstream metric.

The active recovery route is provider-specific: Qwen3.6-35B-A3B benefits most from the diagnosis note ($+4.7$ pp, $p{=}0.036$), while GLM-5.1 ($+4.6$ pp, $p{=}0.037$) and DeepSeek-V4-Pro ($+3.4$ pp) benefit most from the temperature bump.

Appendix~\ref{app:trap_intervention_more} reports the full component statistics and the conservative diagnosis template; Appendix~\ref{app:trap_case_studies} gives case studies showing how the same trigger state can lead to different provider-specific recoveries.
